{{دیوان
| عنوان = خبر تحیر عشق سن نہ جنوں رہا نہ پری رہی
| شاعر = سراج اورنگ آبادی
| ماخذ = ویکی ماخذ
| ماخذ_لنک = https://ur.wikisource.org/wiki/%D8%AE%D8%A8%D8%B1_%D8%AA%D8%AD%DB%8C%D8%B1_%D8%B9%D8%B4%D9%82_%D8%B3%D9%86_%D9%86%DB%81_%D8%AC%D9%86%D9%88%DA%BA_%D8%B1%DB%81%D8%A7_%D9%86%DB%81_%D9%BE%D8%B1%DB%8C_%D8%B1%DB%81%DB%8C
}}

<poem>
خبر تحیر عشق سن نہ جنوں رہا نہ پری رہی
نہ تو تو رہا نہ تو میں رہا جو رہی سو بے خبری رہی

شہ بے خودی نے عطا کیا مجھے اب لباس برہنگی
نہ خرد کی بخیہ گری رہی نہ جنوں کی پردہ دری رہی

چلی سمت غیب سیں کیا ہوا کہ چمن ظہور کا جل گیا
مگر ایک شاخ نہال غم جسے دل کہو سو ہری رہی

نظر تغافل یار کا گلہ کس زباں سیں بیاں کروں
کہ شراب صد قدح آرزو خم دل میں تھی سو بھری رہی

وو عجب گھڑی تھی میں جس گھڑی لیا درس نسخۂ عشق کا
کہ کتاب عقل کی طاق پر جوں دھری تھی تیوں ہی دھری رہی

ترے جوش حیرت حسن کا اثر اس قدر سیں یہاں ہوا
کہ نہ آئینہ میں رہی جلا نہ پری کوں جلوہ گری رہی

کیا خاک آتش عشق نے دل بے نوائے سراجؔ کوں
نہ خطر رہا نہ حذر رہا مگر ایک بے خطری رہی
</poem>
